\documentclass[10pt]{exam}
\printanswers
\usepackage{epsfig,amssymb}
\usepackage{xcolor}
\definecolor{darkred}{rgb}{0.5,0,0}
\definecolor{darkgreen}{rgb}{0,0.5,0}
\usepackage{hyperref}
\usepackage{fullpage}
\usepackage{tikz}
\pagestyle{empty}
\usepackage{listings}
\setlength{\parindent}{0pt}
\setlength{\parskip}{.25cm}
\newcommand{\comment}[1]{}
\boxedpoints
\addpoints
\usepackage{amsmath}
\usepackage{algorithm2e}
\usepackage{url}
\usepackage{enumitem}
\usepackage{graphics}

\pagestyle{headandfoot}
\firstpageheader{\Large Math 309}{{\Large Homework 1}}{\Large Fall 2026}

\begin{document}

\hrule

\vspace{.50cm} Name:
\text{Charlie Pyle}
\vspace{.25cm}
\begin{questions}
\question Use the definition of an even integer to explain why 180 is even.

\begin{solution}
An integer $n$ is even provided that there exists an integer $a$ such that $n = 2a$.

We have $180 = 2(90)$, and $90$ is an integer. Therefore, by the definition of an even integer, $180$ is even.
\end{solution}

\question Identify the hypothesis and conclusion for each of the following conditional statements. One of these statements is not written in ``if... then...'' format. First you must rewrite it in that format.

\begin{parts}
\part If $x$ is a negative real number, then $x^3$ is a negative real number.
\part All cats are cute.
\end{parts}

\begin{solution}
\textbf{Part 1:}

Hypothesis: $x$ is a negative real number.

Conclusion: $x^3$ is a negative real number.

\textbf{Part 2:}\\
The statement ``All cats are cute'' is not written in if-then format. Rewritten in if-then format, it reads: ``If an animal is a cat, then the animal is cute.''

Hypothesis: An animal is a cat.

Conclusion: The animal is cute.
\end{solution}

\question Consider the following statement:

\textbf{Statement:} If $a$ and $b$ are odd integers and $n$ is an even integer, then $a^2n - b^2 + n$ is an odd integer.

\begin{parts}
\part Give two examples supporting this statement.
\part Construct a know-show table outlining an argument you would use to prove the statement above.
\part Write a complete proof for the statement above.
\end{parts}

\begin{solution}
\textbf{Part 1: Two supporting examples}

Let $a = 1$, $b = 1$, and $n = 2$. All three integers satisfy the hypotheses: $a$ and $b$ are odd, and $n$ is even. We compute
\begin{align*}
a^2n - b^2 + n = 1^2(2) - 1^2 + 2 = 2 - 1 + 2 = 3,
\end{align*}
and $3$ is an odd integer, as claimed.

Let $a = 3$, $b = 5$, and $n = 4$. Again $a$ and $b$ are odd, and $n$ is even. We compute
\begin{align*}
a^2n - b^2 + n = 3^2(4) - 5^2 + 4 = 36 - 25 + 4 = 15,
\end{align*}
and $15$ is an odd integer, as claimed.

\textbf{Part 2: Know-show table}

\begin{center}
\begin{tabular}{|c|p{7cm}|p{5cm}|}
\hline
Step & Know & Reason \\ \hline
$P$ & $a$ and $b$ are odd integers, and $n$ is an even integer. & Hypothesis \\ \hline
$P1$ & There exist integers $k$ and $m$ such that $a = 2k+1$ and $b = 2m+1$; there exists an integer $j$ such that $n = 2j$. & Definition of odd integer; definition of even integer \\ \hline
$P2$ & $a^2n - b^2 + n = 2(4k^2j + 4kj + 2j - 2m^2 - 2m - 1) + 1$. & Algebra: substitute $P1$ into $a^2n-b^2+n$ and simplify \\ \hline
$P3$ & $p = 4k^2j + 4kj + 2j - 2m^2 - 2m - 1$ is an integer. & Closure properties of addition and multiplication for integers \\ \hline
\end{tabular}
\end{center}
 
\begin{center}
\begin{tabular}{|c|p{7cm}|p{5cm}|}
\hline
Step & Show & Reason \\ \hline
$Q1$ & There exists an integer $p$ such that $a^2n - b^2 + n = 2p + 1$. & Matches $P3$, since $p$ is an integer \\ \hline
$Q$ & $a^2n - b^2 + n$ is an odd integer. & Definition of odd integer \\ \hline
\end{tabular}
\end{center}

\textbf{Part 3: Complete proof}

\textbf{Theorem.} If $a$ and $b$ are odd integers and $n$ is an even integer, then $a^2n - b^2 + n$ is an odd integer.

\textit{Proof.} We assume that $a$ and $b$ are odd integers and that $n$ is an even integer, and we will prove that $a^2n - b^2 + n$ is an odd integer.

Since $a$ is an odd integer, there exists an integer $k$ such that

\begin{align*}
a = 2k + 1.
\end{align*}

Since $b$ is an odd integer, there exists an integer $m$ such that

\begin{align*}
b = 2m + 1.
\end{align*}

Since $n$ is an even integer, there exists an integer $j$ such that

\begin{align*}
n = 2j.
\end{align*}

Using algebra, we obtain

\begin{align*}
a^2n - b^2 + n &= (2k+1)^2(2j) - (2m+1)^2 + 2j\\
&= (4k^2+4k+1)(2j) - (4m^2+4m+1) + 2j\\
&= 8k^2j + 8kj + 2j - 4m^2 - 4m - 1 + 2j\\
&= 8k^2j + 8kj + 4j - 4m^2 - 4m - 1\\
&= 2(4k^2j + 4kj + 2j - 2m^2 - 2m - 1) + 1.
\end{align*}

An integer $p$ can be defined by

\begin{align*}
p = 4k^2j + 4kj + 2j - 2m^2 - 2m - 1.
\end{align*}

Since $k$, $m$, and $j$ are integers, $p$ is an integer as well, by the closure properties of addition and multiplication for integers. We have shown that

\begin{align*}
a^2n - b^2 + n = 2p + 1,
\end{align*}

where $p$ is an integer. By the definition of an odd integer, we conclude that $a^2n - b^2 + n$ is an odd integer. This completes the proof. $\blacksquare$
\end{solution}

\end{questions}

\end{document}